\documentclass[11pt]{article}
\usepackage[T1]{fontenc}
\usepackage{lmodern}
\usepackage[margin=1.08in]{geometry}
\usepackage{amsmath,amssymb,amsthm,mathtools}
\usepackage{microtype}
\usepackage[numbers,sort&compress]{natbib}
\usepackage[colorlinks=true,linkcolor=blue,citecolor=blue,urlcolor=blue]{hyperref}
\usepackage{enumitem}
\usepackage{booktabs}
\hypersetup{pdftitle={The irrationality exponent of pi is 2},pdfauthor={OpenAI},
  pdfsubject={Rational approximation to pi and convergence of the Flint-Hills series}}
\setlist{topsep=4pt,itemsep=3pt}
\setlength{\emergencystretch}{2em}
\newcommand{\CC}{\mathbb C}
\newcommand{\RR}{\mathbb R}
\newcommand{\QQ}{\mathbb Q}
\newcommand{\ZZ}{\mathbb Z}
\newcommand{\NN}{\mathbb Z_{\ge0}}
\newcommand{\ii}{\mathrm i}
\DeclareMathOperator{\ord}{ord}
\DeclareMathOperator{\lcm}{lcm}
\DeclareMathOperator{\codim}{codim}
\DeclareMathOperator{\Spec}{Spec}
\DeclareMathOperator{\Proj}{Proj}
\DeclareMathOperator{\Supp}{Supp}
\newtheorem{theorem}{Theorem}[section]
\newtheorem{proposition}[theorem]{Proposition}
\newtheorem{lemma}[theorem]{Lemma}
\newtheorem{corollary}[theorem]{Corollary}
\theoremstyle{definition}
\newtheorem{definition}[theorem]{Definition}
\theoremstyle{remark}
\newtheorem{remark}[theorem]{Remark}
\numberwithin{equation}{section}
\title{The irrationality exponent of \(\pi\) is \(2\)}
\author{OpenAI}
\date{September 24, 2026}
\begin{document}
\maketitle

\begin{abstract}
We prove the conjecture that the irrationality exponent of \(\pi\) is \(2\).
As a consequence, the classical Flint--Hills series
\(\sum_{n\ge1}1/(n^3\sin^2 n)\) converges, with angles in radians.
\end{abstract}

\section{Introduction}\label{sec:introduction}

For an irrational real number \(x\), its \emph{irrationality exponent}
(also called its irrationality measure) is
\[
\mu(x)=\sup\left\{\nu>0:
0<\left|x-\frac pq\right|<q^{-\nu}
\text{ for infinitely many coprime }p,q\in\ZZ,\ q\ge2\right\}.
\]
The pigeonhole principle gives \(\mu(x)\ge2\). An upper bound limits
how frequently rational numbers can approximate \(x\) exceptionally
closely; irrationality alone supplies no finite upper bound. The expected
value for the usual positive circle constant is \(2\). Waldschmidt
records the corresponding epsilon-dependent approximation conjecture
in \cite[p.~265]{waldschmidt}. Our main result proves it.

\begin{theorem}\label{thm:main}
The irrationality exponent of \(\pi\) is \(2\). More precisely, for every
real \(\nu>2\), there is an integer \(Q(\nu)\) such that
\[
\left|\pi-\frac pq\right|\ge q^{-\nu}
\qquad(p\in\ZZ,\ q\in\ZZ,\ q\ge Q(\nu)).
\]
\end{theorem}

The inequality applies to all integer numerators and denominators,
whether or not the fraction is reduced. The threshold may depend on
\(\nu\), and the argument does not give it effectively. In particular,
an exponent of \(2\) is different from a uniform positive lower bound
of the form \(c/q^2\), or equivalently from bounded continued-fraction
partial quotients.

The related Flint--Hills problem, popularized by Pickover and discussed
by Alekseyev \cite{alekseyev}, asks whether
\[
\sum_{n\ge1}\frac{1}{n^3\sin^2 n}
\]
converges, with angles in radians. Integers very close to multiples of
\(\pi\) make individual denominators small. A rational-approximation
bound controls how often such small values can occur, not just their
individual size.

\begin{corollary}\label{thm:flinthills}
The classical Flint--Hills series converges.
\end{corollary}

Alekseyev showed that convergence requires
\(\mu(\pi)\le5/2\) \cite[Corollary~4]{alekseyev};
Meiburg proved the sufficient strict condition \(\mu(\pi)<5/2\)
\cite[Theorem~2.5]{meiburg}. Theorem~\ref{thm:main} supplies this
condition. Section~\ref{sec:series} includes a short dyadic proof of the
implication needed here. Neither cited criterion settles its equality
boundary at \(5/2\).

More generally, for fixed real numbers \(a,b>0\), consider the family
\[
\sum_{n=1}^{\infty}\frac{1}{n^a|\sin n|^b}.
\]
Theorem~\ref{thm:main} and the spacing argument in
Section~\ref{sec:series} show that, with angles in radians, this series
converges if and only if \(a>\max\{1,b\}\); see
Corollary~\ref{series:generalized}.

\subsection{Quantitative approximation and related work}

Mahler's 1953 paper \cite{mahler1953} established a finite irrationality
exponent for \(\pi\) by constructing polynomial-coefficient combinations
of powers of logarithms and proving determinant nonvanishing. It gave the explicit
exponent \(42\) for all denominators at least two and also stated an
eventual improvement to \(30\). Mignotte refined that construction,
obtaining exponent \(21\) for all such denominators and \(20\)
eventually \cite[Theorem~1]{mignotte1974}. Chudnovsky developed
Hermite--Pad\'e approximation to exponential functions with precise
asymptotic estimates \cite{chudnovsky1982}.

Hata's integral constructions led to the bound
\(\mu(\pi)\le8.01604539\ldots\) \cite{hata1993}; Zudilin gives
an account of the simultaneous logarithmic approximation and denominator
arithmetic behind this value \cite[Theorem~2]{zudilin2004}.
Salikhov's symmetric-integral method improved it to
\(7.606308\ldots\) \cite{salikhov,salikhov2010}.
Zeilberger and Zudilin combined parameter experiments with a rigorous
arithmetic and asymptotic analysis to obtain
\(7.103205334137\ldots\) \cite{zeilbergerzudilin}.
Bai's September 2026 preprint claims the further bound
\(\mu(\pi)<7.101862832357\) by a two-exponent variation of the latter
integral \cite{Bai2026}. These quantitative bounds are context, not
inputs to our proof.

The connection with trigonometric series also has a longer history:
Mahler applied approximation estimates to sine-denominator Dirichlet
series in \cite[Section~16]{mahler1953}. For the classical Flint--Hills
question, recent work of Dan on generalized series \cite{Dan2026}
and L\'opez Zapata's continued-fraction criterion
\cite{LopezZapata2026} leaves convergence undecided. The earlier
convergence claim of Mantzakouras and L\'opez Zapata was withdrawn
in September 2026 because its proof did not establish convergence
or the asserted irrationality bound \cite{MantzakourasLopezZapata2025}.
Carella claims both exponent two and the stronger bounded-partial-quotient
property \cite{Carella2022}. Neither claim is used here. The displayed
proof of the stronger assertion contains a sign obstruction.\footnote{In the notation of
\cite[Theorem~8.1, Equations~(52) and~(55)]{Carella2022}, let
\(C=p_{n+1}-\pi q_{n+1}-1/q_n\). The convergent estimate
\(\lvert p_{n+1}-\pi q_{n+1}\rvert<1/q_{n+2}<1/q_n\)
gives \(-2/q_n<C<0\). Thus \(\sin C<0\) for large \(n\), contrary
to the positive lower bound in that argument. This objection concerns
the displayed proof of the stronger assertion; it does not decide the
exponent-two claim or the truth of the stronger conclusion.}

\subsection{The proof in outline}

The proof uses a multivariable interpolation theorem to construct a
nonzero determinant, then compares two incompatible estimates for its
absolute value. This organization is related to the interpolation-determinant
method discussed by Laurent \cite[Section~6]{laurent1991}.
Successively separated approximation denominators already play a central
role in Roth's many-variable method \cite[Sections~2--3]{roth1960}.
Our application requires a different interpolation statement, adapted to
logarithmic curves and with thresholds independent of their centers.

\paragraph{Prescribing coefficients.}
A jet here means a finite packet of Taylor coefficients. Give the
coordinates \(Y,X_1,\ldots,X_m\) positive weights and retain polynomial
monomials satisfying one weighted degree inequality. Their exponent
vectors fill a simplex. At each point \((1,c_{j1},\ldots,c_{jm})\),
expand after the substitution
\[
 Y=1+t,\qquad X_i=c_{ji}+u_i+\log(1+t).
\]
The variable \(t\) follows the logarithmic curve; the variables \(u_i\)
are transverse displacements. Theorem~\ref{geom:interpolation} states
that every packet below the prescribed weighted cutoff can be realized,
at finitely many centers whose \(X_i\)-coordinates are pairwise distinct
for each \(i\), simultaneously. The weights of \(X_1,\ldots,X_m\)
are chosen successively before the center locations are known.
The weights leave more polynomial coefficients than prescribed Taylor
coefficients, but that dimension count alone does not prove full rank
at these special centers.

The proof in Section~\ref{geom:section} starts from a local multiplicity
comparison between coordinate directions and commuting differential
directions. Derivative persistence and separated scales restrict any
algebraic obstruction; logarithmic residues then force coordinate
constancy. This stage is related to the zero-estimate methods of
Philippon \cite{philippon1986} and the separated multidegrees in
Farhi's Roth lemma \cite[Section~5.2]{farhi2006}. The resulting curve
inequality bounds total weighted contact at the centers by the weighted
degree of the curve. It gives positivity on an ordinary blowup, and vanishing yields
all the requested coefficients. The relation between curvewise positivity
and asymptotic jet generation is familiar from Demailly's work
\cite[Section~6]{demailly1992}; the weighted, multiple-center statement
on the possibly singular schemes used here is proved explicitly.

\paragraph{Two ways to make the determinant small.}
Suppose that a fixed exponent \(\nu>2\) permits exceptionally close
approximations with unbounded denominators. Choose finitely many of these
successively and set \(c_{ji}=2\ii j p_i/q_i\). They approximate
the common logarithmic periods \(2\pi\ii j\). After truncating the
logarithms by a filtration-preserving change, interpolation gives a
matrix over \(\QQ(\ii)\) with independent rows. A square minor using
all rows is nonzero. Clearing its denominators gives an arithmetic
lower bound.

For the analytic bound, translate each row to the exact periods and
expand in the transverse variables. The index \(a\) denotes the chosen
transverse coefficient. For a polynomial column \(P\),
\[
 [u^a]P(e^z,z+u_1,\ldots,z+u_m)
\]
is an entire function of \(z\). Rows test its Taylor coefficients after
\(z=2\pi\ii j+\log(1+t)\). Rows with the same \(a\) therefore test
the same entire functions. On expanding those functions in \(z\),
repeated Taylor degrees within that group annihilate a
determinant term. Distinct degrees then force a quadratic saving in the
group size. This uses the Taylor-expansion mechanism also appearing in
\cite[Theorem~1]{laurent2000}. If many rows have small weighted
transverse index, there are too few groups to avoid this saving. If
few do, the many large indices instead supply powers of the small
rational-approximation errors. Section~\ref{det:section} proves uniform
bounds for both alternatives before summing the expansion.

\paragraph{Why full simplices and separated choices matter.}
The polynomial degree of \(Y\) and those of the \(X_i\) share one
budget. Both the interpolation volume and the arithmetic column cost
therefore come from full simplices. Comparing the row count in dimension
\(m+1\) with the transverse-index count in dimension \(m\) produces
a collision saving that can grow with \(m\), while the other errors
decrease. Section~\ref{sec:parameters} arranges this for each fixed
\(\nu>2\), then chooses approximation denominators meeting all the
interpolation thresholds. Their independence from centers prevents a
circular choice. Only after dimension, weights, and centers have been
fixed does the polynomial degree tend to infinity. The final section
proves Corollary~\ref{thm:flinthills} using a fixed exponent between
\(2\) and \(5/2\), and establishes the exact convergence criterion in
Corollary~\ref{series:generalized}.

\paragraph{Conventions and inputs.}
All logarithms are natural; \(\log(1+t)\) denotes its power series at
zero. Write \(\ii^2=-1\), and let multi-indices have nonnegative
integer entries. For them, \(\alpha!=\prod_i\alpha_i!\) and
\(\binom\alpha\beta=\prod_i\binom{\alpha_i}{\beta_i}\), with the
latter zero unless \(\beta\le\alpha\) coordinatewise. Geometry is
over \(\CC\), and a curve means an integral algebraic curve. We use
polynomial B\'ezout with isolated local multiplicities, ordinary blowups,
nef-plus-ample, eventual Proj recovery, Serre vanishing, and smooth
projective models of curves. Their exact hypotheses are checked at use;
no prior irrationality bound for \(\pi\) is assumed.

\section{Interpolation with separated weights}
\label{geom:section}

We prove the interpolation statement used in the determinant argument. Its
essential uniformity is that the weights can be chosen before the interpolation
points: each coordinate weight has a lower threshold depending on the earlier
weights, but not on the locations of the points. The degree at which
interpolation becomes surjective is allowed to depend on all the fixed data.

All multi-indices below have nonnegative integer entries. For a positive
rational vector $W=(w_0,\ldots,w_m)$, the $W$-degree of the
monomial $Y^hX_1^{\alpha_1}\cdots X_m^{\alpha_m}$ is
$w_0h+\sum_iw_i\alpha_i$. Write $\mathcal P_W(H)$ for the complex vector
space of polynomials of $W$-degree at most $H$. For a positive rational vector
$V=(v_0,\ldots,v_m)$, let $\mathcal J_V(H)$ be the quotient of
$\CC[[t,u_1,\ldots,u_m]]$ by the ideal of series whose monomials all have
$V$-weight at least $H$. Thus a vector in $\mathcal J_V(H)$ is a packet of
coefficients indexed by
\[
 v_0s+\sum_{i=1}^m v_i\beta_i<H.
\]
The strict inequality in this definition will be used throughout.

\begin{theorem}[Separated-weight interpolation]
\label{geom:interpolation}
Fix integers $m,K\ge1$ and positive rational numbers $w_0,v_0,\theta$ such
that $0<\theta<1$ and
\begin{equation}
\label{geom:volume-hypotheses}
 K\theta^m<1,
 \qquad K\frac{w_0}{v_0}\theta^m<1.
\end{equation}
There are successive lower thresholds for positive rational numbers
$w_1,\ldots,w_m$, depending only on these fixed data and the previously
chosen weights, with the following property. Put
\[
 W=(w_0,w_1,\ldots,w_m),\qquad
 V=(v_0,w_1/\theta,\ldots,w_m/\theta).
\]
For any points
\[
 a_j=(1,c_{j1},\ldots,c_{jm})\in\CC^{m+1},\qquad 0\le j<K,
\]
whose coordinates $c_{0i},\ldots,c_{K-1,i}$ are pairwise distinct for each
$i$, the map
\begin{equation}
\label{geom:jet-map}
 \begin{aligned}
 \mathcal P_W(H)&\longrightarrow\bigoplus_{j=0}^{K-1}\mathcal J_V(H),\\
 P&\longmapsto
 \left(P\bigl(1+t,(c_{ji}+u_i+\log(1+t))_{i=1}^m\bigr)\right)_{j=0}^{K-1}.
 \end{aligned}
\end{equation}
is surjective for all sufficiently large, sufficiently divisible integers $H$.
The divisibility requirement can be chosen using only the weights. The
remaining lower threshold for $H$ may depend on the points.
\end{theorem}

We first establish a numerical inequality on algebraic curves, and then use
it to obtain all the jets in \eqref{geom:jet-map}. The local estimate driving
the curve argument is recorded separately.

\subsection{A local multiplicity comparison}

A collection of vector fields is called a \emph{normal basis} to a smooth
subvariety at a point if its images form a basis of the ambient tangent space
modulo the tangent space of the subvariety. For an irreducible subvariety,
this terminology will refer to a general smooth point. There are only
finitely many possible index sets in the applications below, so all their
generic rank conditions can be checked on one common open subset.

\begin{lemma}[Weighted transverse multiplicity]
\label{geom:multiplicity}
Give the coordinates $x_1,\ldots,x_d$ positive rational degree weights
$\rho_1,\ldots,\rho_d$. On an open subset of $\CC^d$, let
$D_1,\ldots,D_d$ be commuting analytic vector fields forming a frame, and
assign them positive rational costs $\kappa_1,\ldots,\kappa_d$. Fix
$\varepsilon>0$.

Let $N>0$. Suppose that polynomials $f_1,\ldots,f_r$ have weighted degree at most $N$,
that $Z$ is an irreducible component of their common zero set of codimension
$k\ge1$, and that
\begin{equation}
\label{geom:persistent-derivatives}
 D^\gamma f_\ell\big|_Z=0
 \quad\text{whenever}\quad
 \sum_b\kappa_b\gamma_b<\varepsilon N.
\end{equation}
For any coordinate normal basis $A$ and any $D$-field normal basis $B$,
each of cardinality $k$, one has
\begin{equation}
\label{geom:multiplicity-comparison}
 \prod_{a\in A}\rho_a
 \le 2k!\,\varepsilon^{-k}\prod_{b\in B}\kappa_b
\end{equation}
The bound is uniform in the equations, the component $Z$, and the
analytic frame; no lower threshold for $N$ is required.
\end{lemma}

\begin{proof}
Choose a smooth point $x$ of $Z$ where both normal-basis conditions hold and
which lies on no other irreducible component of the common zero set. Flowing
from $Z$ along the fields indexed by $B$ gives local analytic coordinates
$(z,b)$, where $z$ parametrizes $Z$ and $b=(b_1,\ldots,b_k)$ consists of
transverse flow parameters. Commutativity implies that the chosen fields
are the derivatives in the $b$-directions. Thus
\eqref{geom:persistent-derivatives} says that in the expansion
\[
 f_\ell(z,b)=\sum_{\alpha\in\NN^k} f_{\ell,\alpha}(z)b^\alpha,
\]
every coefficient function of $B$-weight less than $\varepsilon N$ vanishes
identically on the chosen neighborhood in $Z$.

Let $S$ be the affine coordinate slice through $x$ obtained by fixing the
coordinates outside $A$. It is transverse to $Z$, and the zero set of the
restricted equations is isolated at $x$. The flow parameters $b$ give
regular analytic coordinates on $S$; the coordinates $z$ become analytic
functions $z(b)$. Substitution cannot generate a missing term of smaller
$B$-weight. Indeed, each surviving summand already contains a monomial
$b^\alpha$ of weight at least $\varepsilon N$, and its analytic coefficient
has only nonnegative powers of $b$. Consequently the ideal of the
restricted equations is contained in the monomial ideal of weight at least
$\varepsilon N$. Its quotient therefore has length at least
\begin{equation}
\label{geom:local-length-lower}
 \#\left\{\alpha\in\NN^k:
       \sum_{b\in B}\kappa_b\alpha_b<\varepsilon N\right\}
 \ge \frac{(\varepsilon N)^k}{2k!\prod_{b\in B}\kappa_b}
\end{equation}
In fact the factor $1/2$ can be omitted: the simplex
$\{x\in\RR_{\ge0}^k:\sum_{b\in B}\kappa_bx_b<\varepsilon N\}$
is covered by the disjoint half-open unit cubes $\alpha+[0,1)^k$
indexed by the admissible multi-indices $\alpha=\lfloor x\rfloor$. Its volume is
$(\varepsilon N)^k/(k!\prod_{b\in B}\kappa_b)$. Thus the displayed
bound holds for every $N>0$. Analytic local coordinates preserve finite
length, so no radius or coefficient bound for the flow map is needed.

In the regular local ring of the $k$-dimensional slice, the ideal generated
by the restricted equations is primary to the maximal ideal. Over the
infinite field $\CC$, $k$ sufficiently general constant linear
combinations of these equations still have an isolated common zero at
$x$: successive generic combinations avoid the finitely many
positive-dimensional components that remain. Their ideal is contained
in the original ideal. Their quotient length is therefore an upper bound
for the length in \eqref{geom:local-length-lower}.

Each of the $k$ combinations has weighted degree at most $N$ in the
variables indexed by $A$. We claim that their local intersection length
is at most
\begin{equation}
\label{geom:weighted-bezout}
 \frac{N^k}{\prod_{a\in A}\rho_a}.
\end{equation}
Choose an integer $M_0$ clearing the denominators of the $\rho_a$, and
make the substitutions
\[
 x_a=\lambda_a+z_a^{M_0\rho_a}\qquad(a\in A).
\]
Choose the constants $\lambda_a$ so that no preimage of the given point
is ramified. There are $\prod_{a\in A}(M_0\rho_a)$ such preimages, each
with the same local intersection length. The substituted polynomials
have ordinary total degree at most $M_0N$. The ordinary B\'ezout bound
for isolated complete-intersection zeros, with their local
multiplicities, bounds the sum of these lengths by $(M_0N)^k$;
see \cite[Corollary~VIII.3, p.~144]{mondal2021}.
The multiplicities there are local quotient lengths, unchanged by
completion \cite[Section~IV.4.2]{mondal2021}.
Dividing gives \eqref{geom:weighted-bezout}.
This application concerns isolated zeros only and permits other
components elsewhere.

Combining \eqref{geom:local-length-lower} and
\eqref{geom:weighted-bezout} proves the result. Neither the
polynomials nor the possibly varying component enters the constant in
\eqref{geom:multiplicity-comparison}.
\end{proof}

\subsection{The curve inequality}

Choose a positive rational number $\sigma$ sufficiently small that
\begin{equation}
\label{geom:sigma}
 \begin{split}
 (1+3\sigma)^{m+1}K(w_0/v_0)\theta^m&<1,\\
 (1+3\sigma)^mK\theta^m&<1,\\
 (1+\sigma)\theta&<1.
 \end{split}
\end{equation}
This choice is possible by \eqref{geom:volume-hypotheses}. Set $v_i=w_i/\theta$
for $i>0$. Fix $L_0=m+2$ and put
\begin{equation}
\label{geom:constant}
 C_* =\max_{1\le k\le m}2k!(L_0/\sigma)^k.
\end{equation}
We choose the positive weights successively so that for every pair of
subsets $A,B\subseteq\{0,\ldots,m\}$ of the same cardinality between
$1$ and $m$,
\begin{equation}
\label{geom:separation}
 \prod_{a\in A}w_a>C_*\prod_{b\in B}v_b
\end{equation}
whenever the largest positive index in their symmetric difference belongs
to $A$.

Here and below a condition involving that largest index is imposed only
when a positive index does differ. The successive choice is possible:
if this largest index is $i$, the ratio in \eqref{geom:separation} is
$w_i$ times a positive expression in the already chosen lower weights
and the fixed zeroth weights. Every higher index occurs in both sets or
in neither; in the former case its ratio is $w_j/v_j=\theta$. Hence the
expression is independent of all higher weights. There are finitely many
pairs of sets. A sufficiently large choice of $w_i$ enforces all
conditions whose largest differing index is $i$, and subsequent choices
do not disturb them. This construction involves no interpolation point.

Let $C$ be an irreducible algebraic curve whose generic point lies in the
affine coordinate chart. Such a curve has a smooth projective model over
$\CC$; see \cite[Theorem~53.2.6 and Lemma~53.2.8, Tag~0BXX]{stacks}.
On that model define
\begin{equation}
\label{geom:curve-degree}
 \deg_W C=
 \sum_P\max\left\{0,-\frac{\operatorname{ord}_P Y}{w_0},
                  -\frac{\operatorname{ord}_P X_1}{w_1},\ldots,
                  -\frac{\operatorname{ord}_P X_m}{w_m}\right\}.
\end{equation}
The order of an identically zero coordinate is understood as $+\infty$;
it contributes no pole. At a branch $P$ through $a_j$, set
\begin{equation}
\label{geom:branch-contact}
 h_P=\min\left\{
 \frac{\operatorname{ord}_P(Y-1)}{v_0},
 \frac{\operatorname{ord}_P(X_1-c_{j1}-\log Y)}{v_1},\ldots,
 \frac{\operatorname{ord}_P(X_m-c_{jm}-\log Y)}{v_m}
 \right\}.
\end{equation}
The logarithm here is its formal expansion at $Y=1$. All entries have
positive order, possibly infinite, and their minimum is finite on a
nonconstant curve.

\begin{proposition}[Curve inequality]
\label{geom:curve-inequality}
With the weights chosen above, every such curve satisfies
\begin{equation}
\label{geom:curve-bound}
 \deg_W C\ge(1+\sigma)\sum_{P\mapsto a_j}h_P,
\end{equation}
where the sum includes all branches through the prescribed points.
\end{proposition}

\begin{proof}
There is nothing to prove if the curve misses all the points. Suppose
instead that \eqref{geom:curve-bound} fails. For every sufficiently large
auxiliary integer $N$, linear algebra gives a nonzero polynomial $F_N$
of $W$-degree at most $N$ whose expansions at all the points have
$V$-order at least $(1+3\sigma)N$. Indeed, the number of its coefficients
has leading term
\[
 \frac{N^{m+1}}{(m+1)!\prod_{a=0}^m w_a},
\]
whereas the number of conditions has leading term
\[
 \frac{K(1+3\sigma)^{m+1}N^{m+1}}
      {(m+1)!\prod_{a=0}^m v_a}.
\]
Their comparison is the first inequality in \eqref{geom:sigma}. This
argument requires no independence among the conditions.

On the open set $Y\ne0$ consider the commuting frame
\begin{equation}
\label{geom:frame}
 D_0=Y\partial_Y+\sum_{i=1}^m\partial_{X_i},\qquad
 D_i=\partial_{X_i}\quad(1\le i\le m).
\end{equation}
It preserves polynomial degree at most $N$. In the formal coordinates of
\eqref{geom:jet-map},
\[
 D_0=(1+t)\partial_t,\qquad D_i=\partial_{u_i}.
\]
Thus differentiation by $D_a$ decreases weighted order by at most $v_a$.
Every $D^\gamma F_N$ with $\sum_av_a\gamma_a\le\sigma N$ retains
order at least $(1+2\sigma)N$ at each point. On a branch $P$, its order
is consequently at least $(1+2\sigma)N h_P$.

A polynomial of $W$-degree at most $N$ restricts to a rational function
on the complete model of $C$ with total pole degree at most
$N\deg_W C$. Since we assumed that \eqref{geom:curve-bound} fails, the
just obtained total zero order exceeds this pole bound. The ordinary
divisor equality for a nonzero rational function therefore forces
\begin{equation}
\label{geom:derivatives-on-curve}
 D^\gamma F_N\big|_C=0
 \qquad\left(\sum_av_a\gamma_a\le\sigma N\right).
\end{equation}
All the assertions about branch orders hold for formal power series:
each monomial has order at least its weight times $h_P$, and
cancellation can only increase the order.

We have converted excess contact into vanishing of many derivatives on
$C$. We next find a component on which that vanishing persists, as in
the derivative-ideal method of zero estimates
\cite[Section~5]{philippon1986}, and compare all its normal bases.
Put $\delta=\sigma N/L_0$. For $0\le r\le L_0$, let $V_r$ be the common
zero set in $Y\ne0$ of the derivatives of cost at most $r\delta$.
Let $C^\circ$ be the part of $C$ in this affine open set. The $V_r$
form a decreasing chain, and \eqref{geom:derivatives-on-curve} puts
$C^\circ$ in $V_{L_0}$. Choose an irreducible component of $V_{L_0}$ containing
$C^\circ$, and extend it backwards to nested irreducible components of the
earlier $V_r$. Every selected component has dimension between $1$ and
$m$: it contains a curve, and it lies in the proper zero set of
$F_N$. Since there are $L_0+1=m+3$ levels, two adjacent selected
components coincide. Denote this persistent component by $Z$, and
write $k=\operatorname{codim}Z$.

Take the derivatives of cost at most $r\delta$ as the equations of the
earlier of these two levels. Any further derivative of cost less than
$\delta$ vanishes along $Z$, because the fields commute and $Z$ is
also a component at the next level. Lemma~\ref{geom:multiplicity}, with
$\varepsilon=\sigma/L_0$, gives
\begin{equation}
\label{geom:all-bases}
 \prod_{a\in A}w_a\le C_*\prod_{b\in B}v_b
\end{equation}
for every coordinate normal basis $A$ and every frame normal basis $B$.
The equations and $Z$ may vary with $N$; the uniformity proved in
Lemma~\ref{geom:multiplicity} is exactly what permits this use.

We show that $Y$ is constant on $Z$, working on the common smooth open
subset where all the normal-basis conditions hold. Suppose first that
some positive coordinate direction is nontangent to $Z$, and choose the
largest such index $i$. Every larger positive index occurs in neither
kind of normal basis, since $D_j=\partial_{X_j}$ for $j>0$.
Some coordinate normal basis $A$ contains $i$. If a frame normal basis
$B$ omitted $i$, the largest positive index where these sets differ
would be $i$. Then
\eqref{geom:separation} would contradict \eqref{geom:all-bases}.
Thus every frame normal basis contains $D_i$.

The other frame vectors consequently fail to span the normal quotient.
Their ambient span is the hyperplane
\[
 \ker(dX_i-dY/Y).
\]
For an ambient hyperplane $Q$ and a tangent space $T$, the image of
$Q$ in the quotient by $T$ fails to be surjective precisely when
$T\subseteq Q$. We obtain
\begin{equation}
\label{geom:log-differential}
 dX_i=dY/Y\quad\text{on }Z.
\end{equation}
This identity forces $Y$ to be constant on the algebraic variety $Z$.
Indeed, if $Y$ were nonconstant, choose a smooth point where $dY\ne0$
and intersect with general algebraic hyperplanes through that point
until a curve remains on which $Y$ is nonconstant. On the smooth
complete model of this curve, \eqref{geom:log-differential} is an
identity of meromorphic differentials. A nonconstant meromorphic
function $Y$ has a zero or pole, where $dY/Y$ has nonzero residue
$\operatorname{ord}_P Y$. An exact differential $dX_i$ has zero
residue at every point. This is a contradiction. It follows also
that $X_i$ is constant, since $dX_i=0$ in characteristic zero.

If instead every positive coordinate direction is tangent, the $m$
independent vectors $\partial_{X_i}$ all lie in the tangent space of
$Z$. Properness forces that tangent space to be their span, and
$dY=0$ on $Z$ again. We conclude in all cases that $Y$ is constant
on $Z$. Since $C^\circ\subseteq Z$ meets a prescribed point with
$Y=1$, we have $Y=1$ on $C^\circ$, and hence on $C$ by density.

It remains to rule out a violating curve in that fiber. We work with
$C^\circ$ in the affine fiber $Y=1$, while computing orders and degrees
on the complete model of $C$ as before. In this fiber, the degree weights
are $w_i$, the order weights are $v_i=w_i/\theta$,
and $h_P=\min_i\operatorname{ord}_P(X_i-c_{ji})/v_i$.
Choose a new auxiliary polynomial in these $m$ variables. The ratio of
the number of vanishing conditions to the number of coefficients now
tends to $K(1+3\sigma)^m\theta^m<1$, by the second inequality
in \eqref{geom:sigma}. Thus, for sufficiently large $N$, there is a
nonzero polynomial $G_N(X_1,\ldots,X_m)$ of weighted degree at most $N$
and order at least $(1+3\sigma)N$ at every center.

An ordinary partial derivative of cost at most $\sigma N$ still has
order at least $(1+2\sigma)N$ at every center and degree at most $N$.
The same zero--pole comparison on $C$ therefore gives
\[
 \partial_X^\gamma G_N\big|_C=0
 \qquad\left(\sum_{i=1}^m v_i\gamma_i\le\sigma N\right).
\]
If $m=1$, this is already impossible: a nonzero polynomial on the affine
line cannot vanish on a curve. Suppose henceforth that $m\ge2$. For
$0\le r\le L_0$, let $U_r$ be the common zero set in this fiber of the
partial derivatives of $G_N$ of cost at most $r\delta$, with
$\delta=\sigma N/L_0$ as before. The curve $C^\circ$ lies in $U_{L_0}$. Choose
nested irreducible components containing $C^\circ$ backwards through this
chain. Their dimensions lie between $1$ and $m-1$, since $G_N$ is
nonzero, so two adjacent components coincide in a component $Z'$.
Its codimension in the fiber satisfies $1\le k\le m-1$.

For the equations at the earlier of these two levels, every further
partial derivative of cost less than $\delta$ vanishes on $Z'$.
Lemma~\ref{geom:multiplicity}, again with $\varepsilon=\sigma/L_0$,
therefore gives \eqref{geom:all-bases} for $Z'$, with the same $C_*$
and with normal-basis index sets contained in $\{1,\ldots,m\}$.

Here the coordinate and frame bases are the same. They must have a
unique index set: two different basis sets could be ordered so that
the largest differing index belongs to $A$, contradicting
\eqref{geom:separation}. The codimension is positive, so this unique
normal basis has some index $i$. The other coordinate directions
fail to span the normal quotient; hence the tangent space is
contained in $\ker dX_i$, and $X_i$ is constant on $Z'$.
Coordinatewise distinctness shows that $Z'$ meets at most one prescribed
point. Every prescribed point of $C$ belongs to $C^\circ\subseteq Z'$,
so the same holds for $C$.

Let that point be $a_j$. Choose a coordinate $X_l$ nonconstant on
$C$. The nonzero rational function $X_l-c_{jl}$ has total pole degree
at most $w_l\deg_W C$, and at each branch through $a_j$ it has order
at least $(w_l/\theta)h_P$. Therefore
\[
 \deg_W C\ge\theta^{-1}\sum_{P\mapsto a_j}h_P
             >(1+\sigma)\sum_{P\mapsto a_j}h_P,
\]
by the last inequality in \eqref{geom:sigma}. This final contradiction
proves \eqref{geom:curve-bound}.
\end{proof}

\subsection{From the curve inequality to all coefficient packets}

We complete the proof of Theorem~\ref{geom:interpolation}. The argument uses
ordinary projective schemes and the ordinary blow-up; smoothness or
normality of the compactification will not be needed. Curve positivity
and asymptotic jet generation are closely related in the Seshadri-constant
approach of \cite[Section~6]{demailly1992}. Here we supply the passage
for our weighted ideals at several points, including the eventual
ordinary-power comparison required on a possibly singular scheme.

Choose an integer $R$ sufficiently large that $R/w_a$ and $R/v_a$ are
positive integers for every index. Embed affine space by all monomials
of $W$-degree at most $R$, including the constant monomial, and let
$X$ be the projective closure of the image. Since these monomials
include each coordinate, the chart of the constant monomial is
precisely the original affine space. Let $L=\mathcal O_X(1)$, a very
ample line bundle.

For any curve meeting this chart,
\begin{equation}
\label{geom:projective-degree}
 \deg(L|_C)=R\deg_W C.
\end{equation}
To see this on its smooth complete model, at each point clear the
largest pole of the homogeneous monomial coordinates. Every monomial
has pole order at most $R$ times the maximum in
\eqref{geom:curve-degree}. A pure power of one coordinate attains
that bound whenever the maximum is positive, and the constant
monomial attains it when the maximum is zero. The resulting sections
have no common zero, proving \eqref{geom:projective-degree}.

To make the local data algebraic, replace $\log(1+t)$ in each
coordinate by a fixed Taylor polynomial $G_i(t)$ whose omitted
terms have $t$-weight strictly greater than $v_i$. Then
\[
 t=Y-1,\qquad u_i'=X_i-c_{ji}-G_i(t)
\]
are regular local coordinates at $a_j$. This replacement preserves
$h_P$. Indeed, if $t$ is not identically zero on a branch, each
omitted term has order divided by $v_i$ strictly larger than
$\operatorname{ord}_P(t)/v_0$; if $t$ is identically zero, its tail
vanishes. It also preserves interpolation surjectivity for every
threshold: the formal coordinate change fixes $t$ and changes
$u_i$ only by terms of weight greater than $v_i$. It and its inverse
preserve the weighted filtration and induce inverse triangular maps
on each finite coefficient packet.

At $a_j$ form the finite-colength ideal
\[
 I_j=\left(t^{R/v_0},(u_1')^{R/v_1},\ldots,(u_m')^{R/v_m}\right).
\]
The points are distinct. These local ideals, together with the unit
ideal away from the points, define a coherent ideal sheaf $I$ on $X$.
For a branch through a center,
\begin{equation}
\label{geom:ideal-valuation}
 \operatorname{ord}_P I
 =\min_a\left\{(R/v_a)\operatorname{ord}_P z_a\right\}
 =Rh_P,
 \qquad z_0=t,\quad z_i=u_i'.
\end{equation}
Let $p:X'\to X$ be the ordinary blow-up of $I$, namely the relative
Proj of its Rees algebra, and write
\[
 I\mathcal O_{X'}=\mathcal O_{X'}(-E).
\]
The scheme $X$ is integral and Noetherian, $I$ is a coherent nonzero
ideal, and $E$ is an effective Cartier divisor. The tautological
bundle $\mathcal O_{X'}(-E)$ is relatively ample; these are the
standard blow-up properties in
\cite[Tags~02OS, 02ND, and~02NS]{stacks}.

Write $A=p^*L$ in additive divisor notation. We claim that
\begin{equation}
\label{geom:nef-class}
 A-(1+\sigma)E\quad\text{is nef}.
\end{equation}
For a noncontracted integral curve $\Gamma$ upstairs whose image
meets the affine chart, the map to its image is birational, since
$p$ is an isomorphism away from the finite center. On the complete
normalizations, \eqref{geom:projective-degree} and
\eqref{geom:ideal-valuation} give
\[
 A\cdot\Gamma=R\deg_W p(\Gamma),\qquad
 E\cdot\Gamma=R\sum_{P\mapsto a_j}h_P.
\]
Proposition~\ref{geom:curve-inequality} gives the desired nonnegative
degree. A curve whose image misses the affine chart avoids the
exceptional divisor. A contracted curve has $A$-degree zero and
positive degree for $-E$ by relative ampleness. These cases account
for every integral curve, proving \eqref{geom:nef-class}.

For sufficiently large $a>1$, the divisor $aA-E$ is ample. One
can see this directly: choose $b$ so that $I\otimes L^b$ is globally
generated. The Rees construction embeds $X'$ in a product
$X\times\mathbb P^r$, with the second hyperplane bundle restricting
to $bA-E$. Tensoring with the first very ample bundle makes
$(b+1)A-E$ ample. Now the identity
\begin{equation}
\label{geom:ample-combination}
 A-E=
 \frac{a-1}{a(1+\sigma)-1}\bigl(A-(1+\sigma)E\bigr)
 +\frac{\sigma}{a(1+\sigma)-1}(aA-E)
\end{equation}
expresses $A-E$ as a positive multiple of a nef class plus a
positive multiple of an ample class. It is therefore ample;
see \cite[Corollary~1.4.10]{lazarsfeld}. This theorem applies to
projective schemes and does not require a smooth blow-up.

For all sufficiently large $n$, the graded-piece and relative
vanishing theorems for the ordinary Rees algebra give
\begin{equation}
\label{geom:rees-pushdown}
 p_*\mathcal O_{X'}(-nE)=I^n,\qquad
 R^q p_*\mathcal O_{X'}(-nE)=0\quad(q>0).
\end{equation}
Indeed, on any affine open of the Noetherian base the Rees algebra
is a Noetherian graded algebra generated in degree one. Its
sufficiently high graded pieces recover the sections of its
associated sheaf on Proj, and the higher cohomology of those twists
vanishes. A finite affine cover gives one common threshold;
see \cite[Tags~0AG6 and~0AG7]{stacks}. In particular,
\eqref{geom:rees-pushdown} uses ordinary powers $I^n$, not their
integral closures, and makes no assertion about small $n$.

By the projection formula, Leray, and Serre vanishing for the ample
line bundle $A-E$,
\begin{equation}
\label{geom:serre-vanishing}
 H^1(X,I^n\otimes L^n)
 =H^1(X',\mathcal O_{X'}(n(A-E)))=0
\end{equation}
for all sufficiently large $n$; see \cite[Tag~0B5U]{stacks}.
The exact sequence for $I^n$ consequently gives a surjection from
$H^0(X,L^n)$ onto all its local coefficient data modulo $I^n$.
The constant projective coordinate trivializes $L$ on the affine
chart. Finite-colength local quotients are unchanged by completion,
so these algebraic classes are also the formal coefficient data used
in the theorem. In these trivializations, every element of $I^n$ has
weighted order at least $nR$, because every displayed generator
of $I$ has weight $R$. Thus
\begin{equation}
\label{geom:ideal-inclusion}
 I^n\subseteq\{\text{local series of weight at least }nR\}.
\end{equation}
Surjectivity modulo $I^n$ therefore implies surjectivity onto the
smaller quotient consisting of the desired packets of weight less
than $nR$. Equality in \eqref{geom:ideal-inclusion} is not needed.

Finally, every section of $L^n$ for sufficiently large $n$ is
supplied by a homogeneous degree-$n$ polynomial in the projective
embedding coordinates. This follows from Serre vanishing applied
to the homogeneous ideal sheaf of $X$ in its projective space.
Such a polynomial restricts on the affine chart to one of
$W$-degree at most $nR$. Taking $H=nR$ and undoing the filtered
Taylor-polynomial coordinate change proves the surjectivity in
\eqref{geom:jet-map}, and completes the proof of
Theorem~\ref{geom:interpolation}.

\section{The arithmetic and analytic determinant estimates}
\label{det:section}

Interpolation now supplies the algebraic rank needed for an
interpolation determinant, in the sense of
\cite[Section~6]{laurent1991}. We estimate one full-row minor from
below arithmetically, and then translate its rows to exact logarithmic
periods to obtain an analytic upper bound.

Throughout this section all parameters except $H$ are fixed.  Write
\[
 w\alpha=\sum_{i=1}^m w_i\alpha_i,
 \qquad w_* =\min_{1\le i\le m}w_i,
 \qquad \omega=2\pi\ii.
\]
Let $\nu>2$, and suppose that integers $p_i$ and $q_i\ge2$ satisfy
\begin{equation}
 \left|\pi-\frac{p_i}{q_i}\right|\le q_i^{-\nu},
 \qquad w_i=\lceil\log q_i\rceil,
 \qquad r_i=\frac{2\ii p_i}{q_i}
 \quad(1\le i\le m).
 \label{det:approximations}
\end{equation}
Assume that $p_i\ne0$ and that $m,K,w_0,v_0,\theta,w_1,\ldots,w_m$
satisfy the hypotheses of Theorem~\ref{geom:interpolation}, with centers
$c_{ji}=jr_i$, $0\le j<K$.  In particular $w_i$ is a positive integer
for $i>0$, while $w_0$ and $v_0$ are positive rational numbers.  Fix
$F_0>2/\theta$, and put
\begin{equation}
 T_i=\left\lceil\frac{F_0w_i}{v_0}\right\rceil,
 \qquad
 G_i(t)=\sum_{1\le k<T_i}\frac{(-1)^{k+1}t^k}{k}.
 \label{det:truncations}
\end{equation}
Empty sums are zero.  The first omitted term has weight
$v_0T_i\ge F_0w_i>w_i/\theta$.  Consequently the substitution
$u_i\mapsto u_i+G_i(t)-\log(1+t)$ is an invertible substitution preserving
the weighted filtration.  It induces an invertible map on every finite
space of coefficients of weight less than $H$.  Thus replacing the
logarithms in Theorem~\ref{geom:interpolation} by these $G_i$ preserves
surjectivity.

\subsection{A full-row-rank minor and its arithmetic size}

Use the monomials
\[
 P(Y,X)=Y^hX^\alpha,
 \qquad h\in\ZZ_{\ge0},\quad \alpha\in\ZZ_{\ge0}^m,
 \qquad w_0h+w\alpha\le H
\]
as columns, and the triples
\[
 \rho=(j,s,\beta),\qquad
 0\le j<K,\quad s\in\ZZ_{\ge0},\quad\beta\in\ZZ_{\ge0}^m,
 \qquad v_0s+\frac{w\beta}{\theta}<H
\]
as rows.  The matrix entry is the coefficient of $t^su^\beta$ in
\begin{equation}
 P\bigl(1+t,jr_1+G_1(t)+u_1,\ldots,jr_m+G_m(t)+u_m\bigr).
 \label{det:matrix}
\end{equation}
For all sufficiently large, sufficiently divisible $H$, choose a square
minor $\Delta_H\ne0$ using every row, as permitted by
Theorem~\ref{geom:interpolation}.  Denote its size by $M=M_H$, and set
\begin{equation}
 \bar b=\bar b_H=\frac{1}{MH}\sum_{\rho=(j,s,\beta)}w\beta.
 \label{det:bbar}
\end{equation}
The elementary lattice count for a simplex gives
\begin{equation}
 M\sim
 \frac{K\theta^mH^{m+1}}{(m+1)!\,v_0\prod_{i=1}^m w_i},
 \qquad 0\le\bar b\le\theta.
 \label{det:row-count}
\end{equation}
Indeed, after scaling by $H$, the lattice boxes of side lengths $H^{-1}$
are Riemann sums for the fixed simplex.  Its volume is
$\theta^m/((m+1)!v_0\prod_iw_i)$; its boundary has volume zero, so the
strict inequality in the row definition does not affect the leading
term.  We always let $H$ tend to infinity only after all the weights and
centers have been fixed.

\begin{lemma}[Arithmetic lower bound]\label{det:arithmetic}
Let $\Lambda=4\log2$, and define
\begin{equation}
 E_{\rm ar}
 =\frac{\Lambda F_0m}{v_0}
  +\Lambda\sum_{i=1}^m\frac1{w_i}+\frac{\theta}{w_*}.
 \label{det:arith-error}
\end{equation}
Then every minor just chosen satisfies
\begin{equation}
 \frac{\log|\Delta_H|}{MH}
 \ge -(1-\bar b)-E_{\rm ar}.
 \label{det:lower}
\end{equation}
\end{lemma}

\begin{proof}
Let $L_i=\operatorname{lcm}(1,\ldots,T_i-1)$, interpreted as $1$ if
$T_i=1$, and put
\[
 D_H=\prod_{i=1}^m L_i^{\lfloor H/w_i\rfloor}.
\]
Scale a column $Y^hX^\alpha$ by $\prod_iq_i^{\alpha_i}$ and a row
$(j,s,\beta)$ by $\prod_iq_i^{-\beta_i}$.  The scaled entry is zero
unless $\alpha\ge\beta$ componentwise; otherwise it is
\[
 \prod_i\binom{\alpha_i}{\beta_i}
 [t^s](1+t)^h
 \prod_i\bigl(q_i(jr_i+G_i(t))\bigr)^{\alpha_i-\beta_i}.
\]
Here $q_ijr_i=2\ii jp_i\in\ZZ[\ii]$.  Since $G_i$ has coefficients in
$L_i^{-1}\ZZ$, the denominator of this entry divides
$\prod_iL_i^{\alpha_i-\beta_i}$, which divides $D_H$.  Thus multiplying
every entry of the scaled matrix by $D_H$ gives a matrix over
$\ZZ[\ii]$ with nonzero determinant.  A nonzero Gaussian integer has
modulus at least one, so
\begin{equation}
 \begin{aligned}
 \log|\Delta_H|
 &\ge -M\log D_H
      -\sum_{\text{columns }\alpha}\sum_i\alpha_i\log q_i\\
 &\quad+\sum_{\text{rows }\beta}\sum_i\beta_i\log q_i.
 \end{aligned}
 \label{det:clearing}
\end{equation}
For each selected column, $\sum_i\alpha_i\log q_i\le w\alpha\le H$.
Also $0\le w_i-\log q_i<1$, and hence
\[
 \sum_{\text{rows }\beta}\sum_i\beta_i\log q_i
 \ge MH\bar b-\sum_{\text{rows }\beta}\sum_i\beta_i
 \ge MH\bar b-\frac{MH\theta}{w_*}.
\]

For completeness, $\log\operatorname{lcm}(1,\ldots,n)\le\Lambda n$.
To see this, the factorial valuation formula shows that each prime
power in $(\ell,2\ell]$ contributes its logarithmic prime weight to
$\log\binom{2\ell}{\ell}$; all the other contributions to the latter
are nonnegative.  Thus the increase of the log least common multiple
across this interval is at most $2\ell\log2$.  Summing over dyadic
intervals up to the smallest power of two above $n$ gives at most
$4n\log2$.  Therefore
\[
 \frac{\log D_H}{H}
 \le\Lambda\sum_i\frac{T_i}{w_i}
 \le\frac{\Lambda F_0m}{v_0}+\Lambda\sum_i\frac1{w_i}.
\]
Substitution in \eqref{det:clearing} proves \eqref{det:lower}.
\end{proof}

\subsection{Translation to entire functions}

For every multi-index $a\in\ZZ_{\ge0}^m$ and every selected column $P$,
define the entire function
\begin{equation}
 f_{a,P}(z)=[u^a]P(e^z,z+u_1,\ldots,z+u_m).
 \label{det:test-functions}
\end{equation}
We regard $(f_{a,P})_P$ as a row vector of entire functions, indexed by
the selected columns.  For $P=Y^hX^\alpha$ it has the explicit entry
\begin{equation}
 f_{a,P}(z)=
 \begin{cases}
 \displaystyle\binom{\alpha}{a}e^{hz}z^{|\alpha|-|a|},&a\le\alpha,\\
 0,&\text{otherwise}.
 \end{cases}
 \label{det:test-monomials}
\end{equation}
In particular only indices $a$ with $wa\le H$ can contribute.

\begin{lemma}[Exact row translation]\label{det:translation}
Every row $\rho=(j,s,\beta)$ of \eqref{det:matrix} is a sum of at most
\begin{equation}
 Q_H=(\lfloor H\rfloor+1)^{2m}
          (\lfloor H/v_0\rfloor+1)
 \label{det:term-count}
\end{equation}
scalar multiples of row vectors of the form
\[
 \bigl([t^\ell]f_{a,P}(j\omega+\log(1+t))\bigr)_P,
 \qquad a\ge\beta,\quad wa\le H,\quad0\le\ell\le s.
\]
Each scalar $\xi$ occurring with index $a$ satisfies
\begin{equation}
 \begin{aligned}
 |\xi|&\le\exp\{-\nu w(a-\beta)+H E_{\rm tr}\},\\
 E_{\rm tr}
 &=\frac{\nu}{F_0}+\frac{\log2}{v_0}
   +\frac{\log4+\log(2K)+\nu}{w_*}.
 \end{aligned}
 \label{det:translation-bound}
\end{equation}
\end{lemma}

\begin{proof}
Put
\[
 \epsilon_i=j(r_i-\omega),\qquad
 \tau_i(t)=G_i(t)-\log(1+t).
\]
Since $e^{j\omega}=1$, the substitutions in \eqref{det:matrix} can be
written as $Y=e^z$, $X_i=z+u_i+\epsilon_i+\tau_i(t)$, where
$z=j\omega+\log(1+t)$.  Expanding first in the variables $u_i$, then
in the constants $\epsilon_i$ and tails $\tau_i$, gives the exact row
vector identity
\begin{align}
 \text{row}_{j,s,\beta}
 ={}&\sum_{\substack{a\ge\beta\\wa\le H}}
       \binom a\beta
       \sum_{0\le d\le a-\beta}\binom{a-\beta}{d}
       \epsilon^{a-\beta-d}
       \sum_{k=0}^s
       \left([t^k]\prod_i\tau_i(t)^{d_i}\right)
 \nonumber\\[-2pt]
 &\hspace{35mm}\cdot
 \bigl([t^{s-k}]f_{a,P}(j\omega+\log(1+t))\bigr)_P.
 \label{det:row-identity}
\end{align}
All scalar factors in this identity are independent of the column
$P$.  This is therefore an identity to which determinant
multilinearity applies directly.

The product of tails vanishes to order at least $\sum_i d_iT_i$.
Consequently a nonzero coefficient in \eqref{det:row-identity}
satisfies
\begin{equation}
 \sum_i d_iw_i\le\frac{v_0}{F_0}\sum_i d_iT_i
 \le\frac{v_0s}{F_0}<\frac{H}{F_0}.
 \label{det:tail-budget}
\end{equation}
On $|t|=1/2$,
\[
 |\tau_i(t)|
 \le\sum_{k\ge T_i}\frac{2^{-k}}k\le1.
\]
Cauchy's coefficient estimate thus bounds the tail coefficient in
\eqref{det:row-identity} by $2^k\le2^s$.  From
\eqref{det:approximations},
\[
 |\epsilon_i|\le2Kq_i^{-\nu}
 \le2K e^\nu e^{-\nu w_i}.
\]
Moreover
$\binom a\beta\binom{a-\beta}{d}\le4^{|a|}$ and
$|a|\le H/w_*$.  Combining these inequalities and using
\eqref{det:tail-budget} gives
\[
 \log|\xi|
 \le-\nu w(a-\beta)+\frac{\nu H}{F_0}
       +\frac{H\log2}{v_0}
       +\frac{H(\log4+\log(2K)+\nu)}{w_*}
\]
for each nonzero scalar.  A zero scalar obeys the required bound as
well.  Finally, each coordinate of $a$ and $d$ is at most $H$ because
$w_i\ge1$, while $0\le k\le H/v_0$.  This proves the term count and
the Lemma.
\end{proof}

\subsection{Repeated Taylor degrees force a quadratic saving}

The Taylor expansion of interpolation determinants eliminates repeated
degrees \cite[Theorem~1]{laurent2000}. Our rows fall into different
transverse-index groups, so this cancellation is used only inside a
fixed group. The next bound is uniform over all row choices.

\begin{lemma}[Collision estimate]\label{det:collision}
Choose one translated test row from each row expansion in
Lemma~\ref{det:translation}, and remove its scalar factor.  Let $T$ be
the resulting square matrix, and let $n_a$ be the number of its rows
with index $a$.  With
\begin{equation}
 c=\frac{\log2}{4},\qquad
 E_{\rm hol}=\frac{100K}{w_0}+\frac{\log2}{v_0}
             +\frac{\log(200K)}{w_*},
 \label{det:hol-error}
\end{equation}
there is a quantity $\rho_H\to0$, independent of all these row
choices and of the selected columns, such that
\begin{equation}
 |\det T|
 \le\exp\left\{-c\sum_a n_a^2+MH(E_{\rm hol}+\rho_H)\right\}.
 \label{det:collision-bound}
\end{equation}
\end{lemma}

\begin{proof}
Put $R=100K$.  Equation~\eqref{det:test-monomials} gives, on $|z|\le R$,
\[
 |f_{a,P}(z)|\le
 \exp\left\{H\left(\frac R{w_0}
                      +\frac{\log(2R)}{w_*}\right)\right\}
 =:\mathcal D_H.
\]
Expand
\[
 f_{a,P}(z)=\sum_{d\ge0}c_{a,d,P}(z/R)^d.
\]
Cauchy's estimate gives $|c_{a,d,P}|\le\mathcal D_H$.  A row of $T$ is
a test of order $\ell\le s$ at some center $j$.  Since
\[
 |j\omega+\log(1+t)|
 \le2\pi(K-1)+\log2<50K=R/2
 \qquad(|t|=1/2),
\]
its value on $(z/R)^d$ has modulus at most $2^\ell2^{-d}$.

Expand the determinant in these power-series rows.  The expansion is
absolutely convergent: a coefficient determinant has modulus at most
$M!\mathcal D_H^M$, and the sum of the product bounds $2^{-d}$ over
all degrees is finite.  If two rows have the same pair $(a,d)$, their
coefficient vectors $(c_{a,d,P})_P$ agree, and that summand vanishes.
Thus every nonzero summand has distinct degrees within each group
$a$, and so
\[
 \sum_{\text{rows}}d\ge\sum_a\binom{n_a}{2}.
\]
Use half of the geometric decay for this lower bound and sum the
other half freely over all nonnegative degrees.  Since
$\sum\ell\le MH/v_0$, this gives
\[
 |\det T|
 \le M!\mathcal D_H^M2^{MH/v_0}
       2^{-\frac12\sum_a\binom{n_a}{2}}
       (1-2^{-1/2})^{-M}.
\]
Because $\sum_a n_a=M$, the asserted inequality follows with, for
example,
\[
 \rho_H=\frac{\log M+\frac14\log2
                    -\log(1-2^{-1/2})}{H}.
\]
Equation~\eqref{det:row-count} shows that $\rho_H\to0$ with the fixed
parameters stipulated above.
\end{proof}

\subsection{A comparison with two alternatives}

Put
\begin{equation}
 E_{\rm an}=E_{\rm tr}+E_{\rm hol}.
 \label{det:analytic-error}
\end{equation}

Let $0<A<1$ set a cutoff $AH$ for transverse weights, and let
$0<\eta<1$ be a fraction of rows. The main arithmetic cost is
$1-\bar b$. In each term of the determinant expansion, if at least
$\eta M$ rows have $wa\le AH$, they give a collision saving;
otherwise the rows with $wa>AH$ supply approximation-error factors.
We seek a saving larger than $1-\bar b+E_{\rm ar}+E_{\rm an}$, and
choose $A$ and $\eta$ in Section~\ref{sec:parameters}.

\begin{proposition}[Determinant comparison]\label{det:comparison}
Let $0<A<1$ and $0<\eta<1$, and define
\begin{equation}
 g=\nu\bigl(A(1-\eta)-\theta\bigr)-(1-\theta),
 \qquad
 L=\frac{\eta^2K\theta^m}{(m+1)v_0A^m}.
 \label{det:budgets}
\end{equation}
The approximations \eqref{det:approximations} and the interpolation
hypotheses cannot hold simultaneously if
\begin{equation}
 g>0,\qquad
 E_{\rm ar}+E_{\rm an}<g,
 \qquad
 cL>1+E_{\rm ar}+E_{\rm an}.
 \label{det:contradiction-conditions}
\end{equation}
\end{proposition}

\begin{proof}
Let
\[
 D_A(H)=\#\{a\in\ZZ_{\ge0}^m:wa\le AH\},
 \qquad L_H=\frac{\eta^2M}{H D_A(H)}.
\]
The simplex count, now in dimension $m$, and
\eqref{det:row-count} show that
\[
 D_A(H)\sim\frac{(AH)^m}{m!\prod_iw_i},
 \qquad L_H\longrightarrow L.
\]
Expand $\Delta_H$ by the row identities in
Lemma~\ref{det:translation}.  For a choice of terms with indices
$a_\rho$, the product of their scalar factors has modulus at most
\[
 \exp\left\{-\nu\sum_\rho wa_\rho
             +\nu MH\bar b+MH E_{\rm tr}\right\}.
\]
There are at most $Q_H^M$ term choices.  Apply
Lemma~\ref{det:collision} to each one.

Suppose first that at least $\eta M$ of its indices obey $wa\le AH$.
There are at most $D_A(H)$ such indices, so Cauchy--Schwarz gives
\[
 \sum_a n_a^2\ge\frac{(\eta M)^2}{D_A(H)}=MH L_H.
\]
The scalar main exponent $-\nu\sum_\rho w(a_\rho-\beta_\rho)$ is
nonpositive, because $a_\rho\ge\beta_\rho$.  Discarding it leaves
an upper bound $-cL_H+E_{\rm an}+\rho_H$ after division by $MH$.

In the other case, more than $(1-\eta)M$ indices have $wa>AH$.
Hence $\sum_\rho wa_\rho\ge MH A(1-\eta)$.  Discarding the
nonpositive collision term leaves the upper bound
$-\nu(A(1-\eta)-\bar b)+E_{\rm an}+\rho_H$.
Taking absolute values before summing the determinant expansion gives
\begin{equation}
 \begin{aligned}
 \frac{\log|\Delta_H|}{MH}
 &\le E_{\rm an}+\varepsilon_H+
       \max\{-cL_H,-\nu(A(1-\eta)-\bar b)\},\\
 \varepsilon_H&=\rho_H+\frac{\log Q_H}{H}\longrightarrow0.
 \end{aligned}
 \label{det:upper}
\end{equation}
All bounds here are uniform over the row-term choices.

For the second alternative, the gap between the two main exponents
is at least $g$, since $\bar b\le\theta$ gives
\[
 \nu(A(1-\eta)-\bar b)-(1-\bar b)
 =\nu A(1-\eta)-1-(\nu-1)\bar b\ge g.
\]
Under \eqref{det:contradiction-conditions}, both quantities inside the
maximum in \eqref{det:upper}, after adding
$E_{\rm an}+\varepsilon_H$, are strictly less than
$-(1-\bar b)-E_{\rm ar}$ for all sufficiently large $H$.  For the
first quantity, use $L_H\to L$ and $1-\bar b\le1$; for the second,
use the displayed uniform gap.  This contradicts
Lemma~\ref{det:arithmetic}.
\end{proof}

\section{Choice of parameters and proof of Theorem~\ref{thm:main}}
\label{sec:parameters}

We now combine the interpolation theorem and the determinant estimates.
The order of choices is essential: the dimension is fixed before the
approximation scales, and all these data are fixed before the polynomial
degree tends to infinity.

\begin{lemma}\label{par:constants}
For every real \(\nu>2\), there are positive rational numbers
\(\theta,A,B,C\) such that
\begin{equation}\label{par:inequalities}
0<\theta<A<B<1,\qquad
\nu(A-\theta)>1-\theta,\qquad
C>1,\quad B<1/C,\quad B<C\theta<1.
\end{equation}
One can then choose \(0<\eta<1\) such that
\begin{equation}\label{par:gap}
g:=\nu\bigl(A(1-\eta)-\theta\bigr)-(1-\theta)>0.
\end{equation}
\end{lemma}

\begin{proof}
Choose a rational \(b\) with \(1/2<b<1-1/\nu\).
For a sufficiently small positive rational \(\delta\), put
\(\theta=1-\delta\) and \(A=1-b\delta\). Then \(0<\theta<A<1\), and
\[
\nu(A-\theta)-(1-\theta)
=\bigl(\nu(1-b)-1\bigr)\delta>0,
\]
whereas
\[
\theta-A^2=(2b-1)\delta-b^2\delta^2>0.
\]
It follows that the interval \((A/\theta,1/A)\) is nonempty.
Choose a rational \(C\) in this interval. Since \(A>\theta\), we have
\(C>1\); since \(\theta<A\), we also have
\(C\theta<\theta/A<1\).
Now choose a rational \(B\) between \(A\) and
\(\min(1/C,C\theta)\). This proves \eqref{par:inequalities}.
Finally, the strict inequality at \(\eta=0\) allows a sufficiently
small \(\eta>0\) in \eqref{par:gap}.
\end{proof}

The two inequalities used in this construction have different roles.
The inequality \(\nu(A-\theta)>1-\theta\) supplies the
approximation-error gap. The inequality \(A^2<\theta\) permits
\(A/\theta<C<1/A\), and hence a choice of \(B\) for which
\(CB<1\), \(C\theta/B>1\), and \(B/A>1\). With the full-simplex
counts of Section~\ref{det:section}, these three ratios make the
dimension-dependent errors tend to zero while the collision saving
tends to infinity. We now choose one finite dimension realizing
both requirements; it will remain fixed throughout the degree limit.

\begin{proof}[Proof of Theorem~\ref{thm:main}]
Fix \(\nu>2\). Suppose for a contradiction that there are arbitrarily
large positive integers \(q\) for which some \(p\in\ZZ\) satisfies
\begin{equation}\label{par:good}
\left|\pi-\frac pq\right|\le q^{-\nu}.
\end{equation}
Fix \(\theta,A,B,C,\eta,g\) from Lemma~\ref{par:constants}, and put
\(\varepsilon_0=\min(g/2,1/2)\).
First choose \(F_0>2/\theta\) so large that
\(\nu/F_0<\varepsilon_0/3\).

For an integer \(m\), set
\begin{equation}\label{par:dimension}
K=\lfloor C^m\rfloor,\qquad
w_0=B^{-m},\qquad
v_0=2K\theta^m w_0.
\end{equation}
These are admissible rational interpolation data:
\begin{equation}\label{par:volume}
K\theta^m\le(C\theta)^m<1,\qquad
K\frac{w_0}{v_0}\theta^m=\frac12.
\end{equation}
The relevant limits as \(m\to\infty\) are
\begin{equation}\label{par:limits}
\frac{K}{w_0}\le(CB)^m\longrightarrow0,\qquad
v_0\sim2\left(\frac{C\theta}{B}\right)^m\longrightarrow\infty,
\qquad
\frac{m}{v_0}\longrightarrow0.
\end{equation}
With the notation of Proposition~\ref{det:comparison}, these choices give
\begin{equation}\label{par:collision-growth}
L=\frac{\eta^2K\theta^m}{(m+1)v_0A^m}
 =\frac{\eta^2(B/A)^m}{2(m+1)}\longrightarrow\infty.
\end{equation}
Thus, with \(\Lambda=4\log2\) and \(c=(\log2)/4\), choose \(m\) large
enough that
\begin{equation}\label{par:dimension-margin}
\frac{\Lambda F_0m+2\log2}{v_0}+\frac{100K}{w_0}
 <\frac{\varepsilon_0}{3},
\qquad cL>2.
\end{equation}
Fix this \(m\), and hence \(K,w_0,v_0\).

The remaining constants are now fixed. Choose approximations
\(p_i/q_i\) satisfying \eqref{par:good}, with
\(w_i=\lceil\log q_i\rceil\), successively sufficiently large for
Theorem~\ref{geom:interpolation} and, in addition, for
\begin{equation}\label{par:weight-margin}
\Lambda\sum_{i=1}^m\frac1{w_i}
+\frac{\theta+\log4+\log(2K)+\nu+\log(200K)}{w_*}
 <\frac{\varepsilon_0}{3},
\qquad w_*=\min_{i>0}w_i.
\end{equation}
This is possible: first impose a sufficiently large lower bound on
every \(w_i\), using \(\sum_i1/w_i\le m/w_*\), and then impose the
successive geometric thresholds. The assumed denominators are
unbounded, so every finite sequence of such requirements can be met.
The thresholds for weight separation are independent of the centers.
For sufficiently large \(q_i\), \(p_i\ne0\), and hence
\(c_{ji}=jr_i\), \(r_i=2\ii p_i/q_i\), are distinct across \(j\) in
each coordinate. No independence between different coordinates is
required.

Using \eqref{det:arith-error}, \eqref{det:translation-bound},
\eqref{det:hol-error}, and \eqref{det:analytic-error}, the sum of errors
is exactly
\[
\begin{aligned}
E_{\rm ar}+E_{\rm an}
&=\frac{\nu}{F_0}
 +\frac{\Lambda F_0m+2\log2}{v_0}+\frac{100K}{w_0}\\
&\quad +\Lambda\sum_i\frac1{w_i}
 +\frac{\theta+\log4+\log(2K)+\nu+\log(200K)}{w_*}.
\end{aligned}
\]
Our three choices make this less than \(\varepsilon_0\).
In particular it is less than \(g\), while
\(cL>2>1+E_{\rm ar}+E_{\rm an}\).
All hypotheses of Proposition~\ref{det:comparison} now hold,
a contradiction. The degree \(H\) in that Proposition tends to infinity
only after the dimension, weights, and centers just chosen are fixed.

We have excluded \eqref{par:good} for arbitrarily large \(q\), for
each fixed \(\nu>2\). This gives the required eventual lower bound,
in fact with a strict inequality. The contradiction allowed unreduced
fractions and zero error. If \(\pi\) were rational, exact representations
with arbitrarily large denominators would contradict it; hence
\(\pi\) is irrational.

Finally divide the unit interval into \(N\) equal parts and apply
pigeonhole to \(0,\pi,\ldots,N\pi\) modulo one. There exist \(p\in\ZZ\)
and \(1\le q\le N\) such that
\[
0<\left|\pi-\frac pq\right|\le\frac1{qN}\le\frac1{q^2}.
\]
After reduction, the same error is at most the inverse square of the
reduced denominator. These errors tend to zero as \(N\to\infty\);
irrationality forces the reduced denominators to be unbounded.
Every exponent less than two therefore occurs infinitely often in
the defining inequality. Hence \(\mu(\pi)\ge2\), and the upper bound
already proved gives \(\mu(\pi)=2\).
\end{proof}

\section{Flint--Hills series and its generalization}\label{sec:series}

We first prove Corollary~\ref{thm:flinthills}. Alekseyev showed
that convergence of the classical Flint--Hills series requires
\(\mu(\pi)\le5/2\) \cite[Corollary~4]{alekseyev}; Meiburg proved
that \(\mu(\pi)<5/2\) is sufficient \cite[Theorem~2.5]{meiburg}.
We include the short spacing argument needed for the consequence of
Theorem~\ref{thm:main}. Write
\(\|x\|=\min_{p\in\ZZ}|x-p|\) for distance to the nearest integer.

\begin{lemma}[Dyadic spacing estimate]\label{series:dyadic}
Let \(\alpha>0\) be irrational. Suppose that, for some \(c>0\) and
\(1<\nu<5/2\),
\[
 \|q\alpha\|\ge c q^{1-\nu}\qquad(q\in\ZZ,\ q\ge1).
\]
Then
\[
 \sum_{q=1}^{\infty}\frac{1}{q^3\|q\alpha\|^2}<\infty.
\]
\end{lemma}

\begin{proof}
Fix a positive integer \(K\), and put
\(d=c(2K)^{1-\nu}\). For \(K\le q<2K\), the points
\(q\alpha\) on the circle \(\RR/\ZZ\) have distance at least
\(d\) from zero. Two distinct such points also have circle distance
at least \(d\): apply the hypothesis to their nonzero integer
difference, whose absolute value is less than \(2K\), and use
\(1-\nu<0\).

On either half-circle starting at zero, arrange the distances in
increasing order. Their first distance is at least \(d\), and each
successive distance increases by at least \(d\). Thus these distances
are at least \(d,2d,3d,\ldots\). No positive integer multiple of
\(\alpha\) lies at either endpoint of a half-circle. Consequently
\[
 \sum_{K\le q<2K}\frac{1}{q^3\|q\alpha\|^2}
 \le \frac{2}{K^3d^2}\sum_{j=1}^{\infty}\frac{1}{j^2}
 =O\bigl(K^{2\nu-5}\bigr).
\]
The implied constant depends only on \(c\) and \(\nu\).
Since \(2\nu-5<0\), summing over \(K=1,2,4,\ldots\) proves
the assertion.
\end{proof}

\begin{proof}[Proof of Corollary~\ref{thm:flinthills}]
Fix a real number \(\nu\) with \(2<\nu<5/2\), and choose a
positive integer \(Q\) as in Theorem~\ref{thm:main}. Taking a nearest
integer numerator in that theorem gives
\[
 \|q\pi\|\ge q^{1-\nu}\qquad(q\ge Q).
\]
The same theorem has established irrationality of \(\pi\). Hence
\[
 c=\min\left(\{1\}\cup
       \left\{q^{\nu-1}\|q\pi\|:q\in\ZZ,\ 1\le q<Q\right\}\right)>0.
\]
It follows that \(\|q\pi\|\ge c q^{1-\nu}\) for every positive
integer \(q\), including the finitely many exceptions below \(Q\).
Lemma~\ref{series:dyadic} therefore gives
\begin{equation}\label{series:distance-sum}
 \sum_{q=1}^{\infty}\frac{1}{q^3\|q\pi\|^2}<\infty.
\end{equation}

For each positive integer \(n\), choose a nearest nonnegative integer
\(q\) to \(n/\pi\), and group the integers \(n\) by \(q\).
Then \(|n-q\pi|\le\pi/2\). The group with
\(q=0\) is finite, and its terms in the Flint--Hills series are finite
because irrationality of \(\pi\) excludes \(\sin n=0\) for a
positive integer \(n\). For each \(q\ge1\), its group lies in an
interval of length \(\pi<4\), so contains at most four integers.
Every integer in that group satisfies
\[
 n\ge\pi(q-1/2)\ge\frac{\pi q}{2}.
\]
Concavity of sine on \([0,\pi/2]\) also gives
\[
 |\sin n|=|\sin(n-q\pi)|
 \ge\frac{2}{\pi}|n-q\pi|
 \ge\frac{2}{\pi}\|q\pi\|.
\]
The contribution of the group indexed by \(q\ge1\) is consequently
at most
\[
 \frac{8}{\pi}\frac{1}{q^3\|q\pi\|^2}.
\]
Summing this bound and using \eqref{series:distance-sum} proves
convergence of \(\sum_{n\ge1}n^{-3}\sin^{-2}n\), with angles in
radians.
\end{proof}


The same spacing argument gives the convergence criterion for the
two-parameter family.

\begin{corollary}[Generalized Flint--Hills series]\label{series:generalized}
For fixed real numbers \(a,b>0\), the series
\[
 \sum_{n=1}^{\infty}\frac{1}{n^a|\sin n|^b}
\]
with angles in radians converges if and only if
\(a>\max\{1,b\}\).
\end{corollary}

\begin{proof}
Suppose first that \(a>\max\{1,b\}\), and choose \(\varepsilon>0\)
so that \(a>\max\{1,b\}+b\varepsilon\).
Theorem~\ref{thm:main}, with the finitely many exceptional denominators
absorbed into a positive constant as above, gives
\(\|q\pi\|\ge c_\varepsilon q^{-1-\varepsilon}\) for all \(q\ge1\).
For an integer \(K\ge1\), put \(d=c_\varepsilon(2K)^{-1-\varepsilon}\).
As in the proof of Lemma~\ref{series:dyadic}, the points \(q\pi\)
on \(\RR/\ZZ\) with \(K\le q<2K\) are mutually \(d\)-separated
and at least \(d\) from zero. There are at most \(K\) points, so the half-circle ordering gives
\[
 \sum_{K\le q<2K}\frac{1}{q^a\|q\pi\|^b}
 \le 2K^{-a}d^{-b}\sum_{j=1}^{K}j^{-b}
 \ll_{a,b,\varepsilon}
 \begin{cases}
 K^{1-a+b\varepsilon},&0<b<1,\\
 K^{1-a+\varepsilon}\log(2K),&b=1,\\
 K^{b-a+b\varepsilon},&b>1.
 \end{cases}
\]
Each power of \(K\) is negative by the choice of \(\varepsilon\);
summing over \(K=1,2,4,\ldots\), including the logarithmic factor
when \(b=1\), proves convergence of
\(\sum_{q\ge1}q^{-a}\|q\pi\|^{-b}\).
Use the same nearest-multiple groups as in the preceding proof. The
\(q=0\) group is finite, while the group indexed by \(q\ge1\)
contributes at most
\[
 4(2/\pi)^a(\pi/2)^b q^{-a}\|q\pi\|^{-b}.
\]
This proves the required convergence.

If \(a\le1\), including the boundary \(a=1\), then
\(|\sin n|\le1\) compares the series from below with the divergent
series \(\sum n^{-a}\). It remains to consider \(1<a\le b\).
Let \(p_j/q_j\) be the continued-fraction convergents to \(\pi\),
whose positive integer numerators tend to infinity. They satisfy
\[
 |\sin p_j|=|\sin(p_j-\pi q_j)|
 \le|p_j-\pi q_j|<q_j^{-1}.
\]
Since \(p_j/q_j\to\pi\), eventually \(p_j\le(\pi+1)q_j\), and hence
\[
 \frac{1}{p_j^a|\sin p_j|^b}
 >p_j^{-a}q_j^b\ge(\pi+1)^{-b}p_j^{b-a}.
\]
These summands are bounded away from zero when \(a=b\), and tend
to infinity when \(a<b\). Thus the series also diverges throughout
\(1<a\le b\), including its boundary \(a=b>1\).
\end{proof}


\renewcommand{\bibfont}{\raggedright}
\bibliographystyle{plainnat}
\bibliography{references}
\end{document}
